\newif\ifglobal

%%%%%%%%%%%%%%%%%%%%%%%
%%% Compile to view %%%
%%%%%%%%%%%%%%%%%%%%%%%

%%%%%%%%%%%%%%%%%%%%%%%%%%%%%%%%%%%%%%%%%%%%%%%%%%%
%%% Readme:                                     %%%
%%% https://www.overleaf.com/read/bwdjvfjksvjy/ %%%
%%%%%%%%%%%%%%%%%%%%%%%%%%%%%%%%%%%%%%%%%%%%%%%%%%%

% >>> M?: marks a module setting number or important notice

% >>> M4: Set {./relative-path} to external files for this module <<<
\def \modulefiles {./01-ULP}
% To add an external file, use:
% (Replace '---' with actual filename)
% '\input{\modulefiles/tables/---.tex}' for tables
% '\includegraphics{\modulefiles/figures/---}' for figures
% '\subfile{\modulefiles/---}' for register files


% >>> M1: This module is in global project? <<<
% 'YES' - uncomment; 'NO' - comment out
\globaltrue

\ifglobal
    \documentclass[main\_\_CN.tex]{subfiles}

\else
    % Font "FandolSong-Regular" used by ctex does not have GSUB table.
    % It causes warnings that do not affect compilation and can be ignored.
    \PassOptionsToPackage{quiet}{fontspec}

    \documentclass[openany, 10pt]{book}

    % >>> M2: In ./00-shared/preamble-trm-module, also find
    %         settings for visibility of todo notes, watermark <<<
    \usepackage{preamble-trm-module}


    % Enable Chinese typesetting
    \usepackage{ctex}

    % >>> M3: Temporary labels in Chinese
    %         you can add your own labels there <<<
    \myexternaldocument{./00-shared/temp-labels-cn}

\fi %\ifglobal


\setcounter{secnumdepth}{4}


% >>> M5: If this module needs any updates in
% './00-shared/preamble-trm-module.sty'
% (include additional package, etc.), contact your module PM
% or create the file './00-shared/changelog.tex' make the
% necessary updates yourself and document them in this file <<<

% Make text left-aligned and allow latex to break pages naturally, comment out for full justification
\raggedright
\raggedbottom

\begin{document}


% Sets language into which pre-defined bits of text are translated
\selectlanguage{chinese}
\setchineseformatting

% Do not delete this block
\ifglobal
% Add the following front matter if
% this module is outside of global project
\else
    \InputIfFileExists{readme}{}{}
    \listoftodos
    {\let\clearpage\relax \listofdonetodos}
    \tableofcontents
    %\listoftables
    %\listoffigures
\fi



%%%%%%%%%%%%%%%%%%%%%%%%%
%%% TRM Module Begins %%%
%%%%%%%%%%%%%%%%%%%%%%%%%

\hypertarget{ulp}{}
\chapter{超低功耗协处理器 (ULP)}
\label{mod:ulp}

\section{概述}

超低功耗协处理器 (ULP, Ultra-Low-Power coprocessor) 是一种功耗极低的处理器设备，可在芯片进入 Deep-sleep 时保持上电（详见低功耗管理章节），允许开发者通过存储在 RTC 存储器中的专用程序，访问 RTC 外围设备、内部传感器及 RTC 寄存器。在对功耗敏感的场景下，主 CPU 处于睡眠状态以降低功耗，协处理器可以由协处理器定时器唤醒，通过控制 RTCIO、RTC I2C、SAR ADC、温度传感器 (TSENS) 等外设检测外部环境或与外部电路进行交互，并在达到唤醒条件时主动唤醒主 CPU。

\begin{figure}[!ht]
    \centering
    \includegraphics[width=1.0\textwidth]{\modulefiles/figures/ulp-overview.pdf}
    \caption{超低功耗协处理器概图}
    \label{figure:coprocessor-overview}
\end{figure}


ESP32-S2 搭载了基于有限状态机 (FSM) 实现的 ULP 协处理器（以下简称 ULP-FSM）和基于 RISC-V 指令集的 ULP 协处理器（以下简称 ULP-RISC-V），用户可根据需求灵活选择。

\section{主要特性}

\begin{itemize}
    \item 可访问最多 8 KB SRAM RTC 慢速内存，储存指令和数据
    \item 采用 8 MHz RTC\_FAST\_CLK 时钟频率
    \item 支持正常模式和 Monitor 模式
    \item 可唤醒 CPU 或向 CPU 发送中断
    \item 可访问外围设备、内部传感器及 RTC 寄存器
\end{itemize}

ULP-FSM 和 ULP-RISC-V 不能同时工作，用户只能选择其中一个作为 ESP32-S2 的超低功耗协处理器。ULP-FSM 与 ULP-RISC-V 特性比较如下：


\begin{table}[!ht]
    \centering
    \begin{tabular}[c]{| p{2.5cm} | p{3cm} | c | c |}
    \hline
    \multicolumn{2}{|c}{} & \multicolumn{2}{|c|}{超低功耗协处理器}         \\
    \cline{3-4}
    \multicolumn{2}{|c|}{\multirow{-2}{*}{特性}} & ULP-FSM & ULP-RISC-V  \\
    \hline
    \multicolumn{2}{|c}{内存（RTC 慢速内存）}  & \multicolumn{2}{|c|}{8 KB}  \\
    \hline
    \multicolumn{2}{|c}{工作时钟频率} & \multicolumn{2}{|c|}{8 MHz} \\
    \hline
    \multicolumn{2}{|c}{唤醒源} & \multicolumn{2}{|c|}{ULP 定时器} \\
    \hline
    \multirow{2}{*}{工作模式} & 正常模式 & \multicolumn{2}{c|}{当系统唤醒时，协助主 CPU 完成部分任务} \\
    \cline{2-4}
    & Monitor 模式 & \multicolumn{2}{c|}{当系统休眠时，可以通过控制传感器完成监控外部环境等任务} \\
    \hline
    \multicolumn{2}{|c}{\multirow{5}{*}{可控制的低功耗外设}} & \multicolumn{2}{|c|}{ADC1/ADC2}    \\
    \cline{3-4}
     \multicolumn{2}{|c}{} & \multicolumn{2}{|c|}{DAC1/DAC2}  \\
    \cline{3-4}
    \multicolumn{2}{|c}{} & \multicolumn{2}{|c|}{RTC I2C}  \\
    \cline{3-4}
    \multicolumn{2}{|c}{} & \multicolumn{2}{|c|}{RTC GPIO}  \\
     \cline{3-4}
    \multicolumn{2}{|c}{} & \multicolumn{2}{|c|}{触摸传感器}  \\
     \cline{3-4}
    \multicolumn{2}{|c}{} & \multicolumn{2}{|c|}{温度传感器}   \\
    \hline
    \multicolumn{2}{|c|}{架构} & 可编程有限状态机 & RISC-V\\
    \hline
    \multicolumn{2}{|c|}{开发} & 特殊指令 & 标准 C 编译器 \\
    \hline
    \end{tabular}
    \caption{超低功耗协处理器特性比较}
\end{table}

ESP32-S2 协处理器的功能灵活，可以通过 RTC 寄存器控制 RTC 域中的模块。协处理器可独立于 CPU 运行，是 CPU 的有力补充，甚至可以在一些功耗敏感的设计中取代 CPU。ESP32-S2 协处理器的基本架构可见图 \ref{figure:coprocessor}。

\begin{figure}[H]
    \centering
    \includegraphics[width=0.7\textwidth]{\modulefiles/figures/co-processor.eps}
    \caption{超低功耗协处理器基本架构}
    \label{figure:coprocessor}
\end{figure}



\section{编程流程}

ULP-RISC-V 支持用户使用 C 语言编写程序，然后使用编译器将程序编译成 \href{https://riscv.org/specifications/}{RV32IMC} 标准指令码。ULP-RISC-V 支持 RV32IMC 指令集。ULP-FSM 不支持高级语言，用户需使用 ULP-FSM 专门指令集进行编程，见章节 \ref{section:instruction-set}。

\begin{figure}[!ht ]
    \centering
    \includegraphics[width=0.7\textwidth]{\modulefiles/figures/ulp-instruction.eps}
    \caption{编程流程图}
    \label{figure:programming-workflow}
\end{figure}

\section{协处理器的睡眠和唤醒流程}

ESP32-S2 的协处理器经过专门设计，无论 CPU 是否处于休眠状态，均可独立于 CPU 运行。

在典型场景中，为了降低功耗，系统可进入 Deep-sleep 模式。系统进入睡眠模式前需完成以下操作：

\begin{enumerate}
    \item 将协处理器需要执行的程序烧写到 RTC 慢速内存；
    \item 配置 \hyperref[fielddesc:RTCCNTLCOCPUSEL]{RTC\_CNTL\_COCPU\_SEL} 寄存器，选择协处理器；
    \begin{itemize}
        \item 0：选择使用 ULP-RISC-V
        \item 1：选择使用 ULP-FSM
    \end{itemize}
    \item 如果选择使用 ULP-RISC-V，则还需置位 \hyperref[fielddesc:RTCCNTLCOCPUCLKFO]{RTC\_CNTL\_COCPU\_CLK\_FO}；
    \item 配置 \hyperref[regdesc:RTCCNTLULPCPTIMER1REG]{RTC\_CNTL\_ULP\_CP\_TIMER\_1\_REG} 寄存器来设置硬件定时器的唤醒间隔时间；
    \item 选择定时器使能方式：
        \begin{itemize}
            \item 软件使能：软件置位 \hyperref[fielddesc:RTCCNTLULPCPSLPTIMEREN]{RTC\_CNTL\_ULP\_CP\_SLP\_TIMER\_EN}；
        \item RTC GPIO 使能：软件配置 \hyperref[fielddesc:RTCCNTLULPCPGPIOWAKEUPENA]{RTC\_CNTL\_ULP\_CP\_GPIO\_WAKEUP\_ENA} 开启通过 RTC GPIO 使能硬件定时器选项。更多信息见章节 \ref{mod:lowpowm} \textit{\nameref{mod:lowpowm}}。
        \end{itemize}
    \item 配置系统进入睡眠状态，主 CPU 休眠。
\end{enumerate}

在 Deep-sleep 模式下：

\begin{enumerate}
   \item 硬件定时器周期性地将低功耗控制器置于 Monitor 状态，然后唤醒协处理器；
   \item 协处理器唤醒后执行一些必要操作，例如通过低功耗传感器监控芯片外部环境等操作。
   \item 操作完成后，系统返回 Deep-sleep 模式；
   \item 协处理器进入休眠，等待下一次唤醒。
\end{enumerate}

进入 Monitor 模式后，协处理器的唤醒和睡眠流程见图 \ref{figure:ULP-program-control-flow}，包括：

\begin{figure}[!ht]
    \centering
    \includegraphics[width=0.9\textwidth]{\modulefiles/figures/ulp-exec-control-flow.eps}
    \caption{ULP 程序流程框图}
    \label{figure:ULP-program-control-flow}
\end{figure}

\begin{enumerate}
    \item 使能硬件定时器，定时器开始计数；
    \item 硬件定时器过期，唤醒协处理器。协处理器进入Run状态，执行预先烧写的程序；
    \item 协处理器执行 HALT 相关操作进入 HALT 状态；协处理器程序停止运行开始休眠。定时器再次启动，重复以上流程；
    \begin{itemize}
        \item ULP-RISC-V 的 HALT 操作：置位寄存器 \hyperref[fielddesc:RTCCNTLCOCPUDONE]{RTC\_CNTL\_COCPU\_DONE}；
        \item ULP-FSM 的 HALT 操作：执行 HALT 指令。
    \end{itemize}
    \item 通过协处理器程序或软件关闭硬件定时器，系统将不再进入Monitor状态。
        \begin{itemize}
            \item 软件关闭：软件清零 \hyperref[fielddesc:RTCCNTLULPCPSLPTIMEREN]{RTC\_CNTL\_ULP\_CP\_SLP\_TIMER\_EN}；
            \item RTC GPIO 关闭：软件清零 \hyperref[fielddesc:RTCCNTLULPCPGPIOWAKEUPENA]{RTC\_CNTL\_ULP\_CP\_GPIO\_WAKEUP\_ENA}，并置位 \hyperref[fielddesc:RTCCNTLULPCPGPIOWAKEUPCLR]{RTC\_CNTL\_ULP\_CP\_GPIO\_WAKEUP\_CLR}。

        注：硬件定时器的关闭方法需要与使能方法一致。
        \end{itemize}
\end{enumerate}

\textbf{注意：}\\
    由于 ULP-RISC-V 进入 HALT 状态需要配置寄存器 \hyperref[fielddesc:RTCCNTLCOCPUDONE]{RTC\_CNTL\_COCPU\_DONE}，所以建议预先烧写的程序用下列代码结尾：
    \begin{itemize}
        \item 置位寄存器 \hyperref[fielddesc:RTCCNTLCOCPUDONE]{RTC\_CNTL\_COCPU\_DONE}，结束本次 ULP-RISC-V 的运行，并使 ULP-RISC-V 进入休眠模式；
        \item 置位寄存器 \hyperref[fielddesc:RTCCNTLCOCPUSHUTRESETEN]{RTC\_CNTL\_COCPU\_SHUT\_RESET\_EN}，复位 ULP-RISC-V。
        注：硬件已经预留了足够的时间支持 ULP-RISC-V 在进入休眠之前完成该操作。
    \end{itemize}

上述信号与寄存器之间的关系可见图 \ref {figure:ULP-program-control}。\\


\begin{figure}[H]
    \centering
    \includegraphics[width=0.8\textwidth]{\modulefiles/figures/ulp-exec-control.eps}
    \caption{ULP 程序框图}
    \label{figure:ULP-program-control}
\end{figure}

\section{ULP-FSM}
\subsection{ULP-FSM 主要特性}

ULP-FSM 协处理器是一种可编程有限状态机，可在 CPU 进入 Deep-sleep 状态时工作。协处理器支持部分通用 CPU 指令，可进行一些复杂的逻辑与算术运算。此外，ULP-FSM 还支持一些特殊的 RTC 控制与外围设备控制指令。ULP-FSM 的运行代码和数据存储在 8 KB SRAM RTC 慢速内存中（CPU 也可访问该区域）。因此，这块内存经常用于存储一些协处理器和 CPU 的通用指令。ULP-FSM 可通过执行 HALT 指令停止运行。

ULP-FSM 具有以下特性：
\begin{itemize}
    \item 采用 4 个 16 位通用寄存器 (R0 \verb+~+ R3)，进行数据操作和内存访问。
    \item 采用 1 个 8 位阶段计数器寄存器 Stage\_cnt，可通过 ALU 指令进行操作并用于 JUMP 指令。
    \item 内置专用指令，可直接控制低功耗外设，如 SAR ADC，温度传感器等。
\end{itemize}

\subsection{ULP-FSM 指令集} \label{section:instruction-set}

ULP-FSM 可支持下列指令：

\begin{itemize}
    \item ALU - 算数与逻辑
    \item LD、ST、REG\_RD 及 REG\_WR - 加载与数据存储
    \item JUMP - 跳转至某地址
    \item WAIT 和 HALT - 管理程序执行
    \item WAKE - 唤醒 CPU 及与 CPU 通信
    \item TSENS 和 ADC - 测量
\end{itemize}

ULP-FSM 指令的格式见图 \ref{figure:ULP-instruction}。

\regfield{OpCode}{4}{28}{{}}%
\regfield{Operands}{28}{0}{{}}%
\begin{figure*}[!ht]
    \centering
    \caption{ULP-FSM 协处理器的指令格式}
    \label{figure:ULP-instruction}
\end{figure*}

根据 \textit{Operands} 的设置不同，同一个 \textit{OpCode} 可对应多种不同操作。例如，ALU 能够执行 10 种不同的算数和逻辑运算，JUMP 也可执行有条件跳转、无条件跳转、绝对跳转及相对跳转等多种形式的跳转。

ULP-FSM 协处理器的所有指令均固定为 32 位。通过这一系列指令，协处理器程序即可得到执行。程序内部的执行均采用 32 位寻址。该程序具体存储在 1 块专用的慢速内存区，地址范围为 \hexprefix{}5000\_0000 到 \hexprefix{}5000\_1FFF (8 KB)，对主 CPU 可见。



\subsubsection{ALU - 算数与逻辑运算} \label{paragraph:ALU}

算数逻辑单元 (ALU) 可以进行算数和逻辑运算，操作对象为协处理器寄存器中存储的数值或指令中存储的立即值。具体可以支持的运算类型如下：
\begin{itemize}
    \item 算数 - 加 (ADD) 和减 (SUB)
    \item 逻辑 - 按位与 (AND) 和按位或 (OR)
    \item 移位 - 左移 (LSH) 和右移 (RSH)
    \item 寄存器赋值 - 移动 (MOVE)
    \item 计数器寄存器操作 - STAGE\_RST、STAGE\_INC 和 STAGE\_DEC
\end{itemize}

尽管此处 \textit{OpCode} 相同，均为 7，但可通过设置协处理器指令 [27:21] 位，选择特定的算数和逻辑运算。

\textcolor{black}{\bfseries 对寄存器数值的运算}

\regfield{}{4}{28}{{7}}%
\regfield{}{2}{26}{{0}}%
\regfield{}{1}{25}{{}}%
\regfield{ALU\_sel}{4}{21}{{}}%
\regfield{}{15}{6}{{}}%
\regfield{Rsrc2}{2}{4}{{}}%
\regfield{Rsrc1}{2}{2}{{}}%
\regfield{Rdst}{2}{0}{{}}%
\begin{figure*}[!ht]
    \centering
    \caption{指令类型 - 对寄存器数值的 ALU 运算}
    \label{figure:ALU-among-registers}
\end{figure*}

如图 \ref{figure:ALU-among-registers} 所示，当指令 [27:26] 位设置为 0 时，ALU 将对寄存器 R[0-3] 中存储的内容进行运算，运算类型则取决于指令的 ALU\_sel [24:21] 位，具体设置方式见表 \ref{table:ALU-among-registers}。

\begin{tabular}{ l l }
\textbf{Operand} & \textbf{描述} - 见图 \ref{figure:ALU-among-registers} \\
\textit{Rdst} & 寄存器 R[0-3]，目标地址寄存器\\
\textit{Rsrc1} & 寄存器 R[0-3]，源地址寄存器 \\
\textit{Rsrc2} & 寄存器 R[0-3]，源地址寄存器 \\
\textit{ALU\_sel} & ALU 运算类型，具体见表 \ref{table:ALU-among-registers}\\
\end{tabular}

\begin{table}[!ht]
    \centering
    \begin{tabular}{ c l l l }
    \textbf{ALU\_sel} & \textbf{指令} & \textbf{运算类型} & \textbf{描述} \\ \hline
    0 & ADD & \textit{Rdst} = \textit{Rsrc1} + \textit{Rsrc2} & 加 \\
    1 & SUB & \textit{Rdst} = \textit{Rsrc1} - \textit{Rsrc2} & 减 \\
    2 & AND & \textit{Rdst} = \textit{Rsrc1} \& \textit{Rsrc2} & 按位与 \\
    3 & OR & \textit{Rdst} = \textit{Rsrc1} | \textit{Rsrc2} & 按位或 \\
    4 & MOVE & \textit{Rdst} = \textit{Rsrc1} & 寄存器赋值 \\
    5 & LSH & \textit{Rdst} = \textit{Rsrc1} \textless\textless\textvisiblespace \textit{Rsrc2} & 左移 \\
    6 & RSH & \textit{Rdst} = \textit{Rsrc1} \textgreater\textgreater\textvisiblespace \textit{Rsrc2} & 右移 \\
    \end{tabular}
    \caption{对寄存器数值的 ALU 运算}
    \label{table:ALU-among-registers}
\end{table}

\newpage
\textbf{注意：}
\begin{itemize}
    \item ADD 和 SUB 运算可用于设置或清除 ALU 溢出标志位。
    \item 所有 ALU 运算均可用于设置或清除 ALU 零标志位。
\end{itemize}

\textcolor{black}{\bfseries 对指令立即值的运算}

\regfield{}{4}{28}{{7}}%
\regfield{}{2}{26}{{1}}%
\regfield{}{1}{25}{{}}%
\regfield{ALU\_sel}{4}{21}{{}}%
\regfield{}{1}{20}{{}}%
\regfield{Imm}{16}{4}{{}}%
\regfield{Rsrc1}{2}{2}{{}}%
\regfield{Rdst}{2}{0}{{}}%
\begin{figure*}[!ht]
    \centering
    \caption{指令类型 - 对指令立即值的 ALU 运算}
    \label{figure:ALU-with-immediate-value}
\end{figure*}

如图 \ref{figure:ALU-with-immediate-value} 所示，当指令 [27:26] 位设置为 1 时，ALU 将对寄存器 R[0-3] 和指令 [19:4] 位存储的立即值进行运算。运算类型取决于指令的 ALU\_sel [24:21] 位，具体设置方式见表 \ref{table:ALU-with-immediate-value}。

\begin{tabular}{ l l }
\textbf{Operand} & \textbf{描述} - 见图 \ref{figure:ALU-with-immediate-value} \\
\textit{Rdst} & 寄存器 R[0-3]，目标地址寄存器 \\
\textit{Rsrc1} & 寄存器 R[0-3]，源地址寄存器 \\
\textit{Imm} & 指令立即值，16 位有符号数 \\
\textit{ALU\_sel} & ALU 运算类型，具体见表 \ref{table:ALU-with-immediate-value}
\end{tabular}

\begin{table}[!ht]
    \centering
    \begin{tabular}{ c l l l }
    \textbf{ALU\_sel} & \textbf{指令} & \textbf{运算} & \textbf{描述} \\ \hline
    0 & ADD & \textit{Rdst} = \textit{Rsrc1} + \textit{Imm} & 加 \\
    1 & SUB & \textit{Rdst} = \textit{Rsrc1} - \textit{Imm} & 减 \\
    2 & AND & \textit{Rdst} = \textit{Rsrc1} \& \textit{Imm} & 按位与 \\
    3 & OR & \textit{Rdst} = \textit{Rsrc1} | \textit{Imm} & 按位或 \\
    4 & MOVE & \textit{Rdst} = \textit{Imm} & 寄存器赋值 \\
    5 & LSH & \textit{Rdst} = \textit{Rsrc1} \textless\textless\textvisiblespace \textit{Imm} & 左移 \\
    6 & RSH & \textit{Rdst} = \textit{Rsrc1} \textgreater\textgreater\textvisiblespace \textit{Imm} & 右移 \\
    \end{tabular}
    \caption{对指令立即值的 ALU 运算}
    \label{table:ALU-with-immediate-value}
\end{table}

\textbf{注意：}
\begin{itemize}
    \item ADD 和 SUB 运算可用于设置或清除 ALU 溢出标志位。
    \item 所有 ALU 运算均可用于设置或清除 ALU 零标志位。
\end{itemize}

\newpage

\textcolor{black}{\bfseries 对阶段计数器寄存器数值的运算}\label{section:alu-cnt}


\regfield{}{4}{28}{{7}}%
\regfield{}{2}{26}{{2}}%
\regfield{}{1}{25}{{}}%
\regfield{ALU\_sel}{4}{21}{{}}%
\regfield{}{9}{12}{{}}%
\regfield{Imm}{8}{4}{{}}%
\regfield{}{4}{0}{{}}%
\begin{figure*}[!ht]
    \centering
    \caption{指令类型 - 对阶段计数器寄存器的 ALU 运算}
    \label{figure:ALU-with-stage-count}
\end{figure*}

如图 \ref{figure:ALU-with-stage-count} 所示，当指令 [27:26] 位设置为 2 时，ALU 将对 8 位寄存器 Stage\_cnt 进行递增、递减或重置操作，运算类型取决于指令的 ALU\_sel [24:21] 位，具体设置方式见表 \ref{figure:ALU-with-stage-count}。Stage\_cnt 是一个独立的寄存器，并不是图 \ref{figure:ALU-with-stage-count} 所示指令的一部分。

\begin{tabular}{ l l }
\textbf{Operand} & \textbf{描述} - 见图 \ref{figure:ALU-with-stage-count} \\
\textit{Imm} & 指令立即值，8 位数 \\
\textit{ALU\_sel} & ALU 运算类型，具体见表 \ref{figure:ALU-with-stage-count}\\
\textit{Stage\_cnt} & 专用 8 位阶段计数器寄存器，可存储循环下标等变量
\end{tabular}


\begin{table}[!ht]
    \centering
    \begin{tabular}{ c l l l }
    \textbf{ALU\_sel} & \textbf{指令} & \textbf{运算} & \textbf{描述} \\ \hline
    0 & STAGE\_INC & \textit{Stage\_cnt} = \textit{Stage\_cnt} + \textit{Imm} & 阶段计数器寄存器递增 \\
    1 & STAGE\_DEC & \textit{Stage\_cnt} = \textit{Stage\_cnt} - \textit{Imm} & 阶段计数器寄存器递减 \\
    2 & STAGE\_RST & \textit{Stage\_cnt} = 0 & 阶段计数器寄存器复位
    \end{tabular}
    \caption{对阶段计数器寄存器的 ALU 运算}
    \label{table:ALU-with-stage-count}
\end{table}

\textbf{注意：} \\
    该指令主要是与基于阶段计数器的 JUMPS 指令配合使用，构成基于阶段计数器的 for 循环。
    用法可参考以下伪代码：

    \begin{table}[!ht ]
    \centering
    \begin{tabular}{ l l }
    \textit{STAGE\_RST}  & // 清空阶段计数器   \\
    \textit{STAGE\_INC} & // 阶段计数器 ++  \\
    \textit{\{...\}} &  // 循环主体，包含 \begin{math}n\end{math} 条指令   \\
    \textit{JUMPS (step = \begin{math}n\end{math}, cond = 0, threshold = \begin{math}m\end{math})}  & // 如果阶段计数值小于 \begin{math}m\end{math}，则跳转至 STAGE\_INC，否则跳出 \\
    \multicolumn{2}{l}{该循环，以此来实现一个阈值为 \begin{math}m\end{math} 的累加 for 循环。}
    \end{tabular}
\end{table}

\subsubsection {ST - 存储数据至内存}\label{st instruction}


\regfield{}{4}{28}{6}%
\regfield{manul\_en}{1}{27}{{}}%
\regfield{offset\_set}{1}{26}{{}}%
\regfield{wr\_auto}{1}{25}{{}}%
\regfield{}{4}{21}{{}}%
\regfield{offset}{11}{10}{{}}%
\regfield{}{1}{9}{{}}%
\regfield{wr\_way}{2}{7}{{}}%
\regfield{upper}{1}{6}{{}}%
\regfield{label}{2}{4}{{}}%
\regfield{Rsrc}{2}{2}{{}}%
\regfield{Rdst}{2}{0}{{}}%
\begin{figure*}[!ht]
    \centering
    \caption{指令类型 - ST}
    \label{figure:ST}
\end{figure*}

\begin{tabular}{ l l }
\textbf{Operand} & \textbf{描述} - 见图 \ref{figure:ST} \\
\textit{Rdst} & 寄存器 R[0-3]，目标寄存器地址，单位为 32 位字\\
\textit{Rsrc} & 寄存器 R[0-3]，保留需存储的 16 位数 \\
\textit{label} & 数据标志，用户自定义的 2 位无符号数\\
\textit{upper} & 0：写低半字；1：写高半字 \\
\textit{wr\_way} & 0：写全字；1：带 label；3：不带 label \\
\textit{offset} & 11 位有符号数，单位为 32 位字\\
\textit{wr\_auto} & 使能地址自增模式 \\
\textit{offset\_set} & offset 使能位。0：不设置地址自增模式的基地址偏移；1：设置地址自增模式下的基地址偏移。\\
\textit{manul\_en} & 使能地址指定模式
\end{tabular}

\textcolor{black}{\bfseries 地址自增模式}

\regfield{}{4}{28}{6}%
\regfield{}{3}{25}{{3}}%
\regfield{}{4}{21}{{}}%
\regfield{offset}{11}{10}{{}}%
\regfield{}{10}{0}{{}}%
\begin{figure*}[!ht]
    \centering
    \caption{指令类型 - 地址自增模式的基地址偏移 (ST-OFFSET)}
    \label{figure:ST-AUTO-OFFSET}
\end{figure*}

\begin{tabular}{ l l }
\textbf{Operand} & \textbf{描述} - 见图 \ref{figure:ST-AUTO-OFFSET} \\
\textit{offset} & 初始地址偏移量，11 位有符号数，单位为 32 位字。 \\
\end{tabular}

\regfield{}{4}{28}{6}%
\regfield{}{3}{25}{{1}}%
\regfield{}{16}{9}{{}}%
\regfield{wr\_way}{2}{7}{{}}%
\regfield{}{1}{6}{{}}%
\regfield{label}{2}{4}{{}}%
\regfield{Rsrc}{2}{2}{{}}%
\regfield{Rdst}{2}{0}{{}}%
\begin{figure*}[!ht]
    \centering
    \caption{指令类型 - 地址自增模式的数据存储 (ST-AUTO-DATA)}
    \label{figure:ST-AUTO-DATA}
\end{figure*}

\begin{tabular}{ l l }
\textbf{Operand} & \textbf{描述} - 见图 \ref{figure:ST-AUTO-DATA} \\
\textit{Rdst} & 寄存器 R[0-3]，目标寄存器地址，单位为 32 位字 \\
\textit{Rsrc} & 寄存器 R[0-3]，保留需存储的 16 位数 \\
\textit{label} & 数据标志，用户自定义的 2 位无符号数 \\
\textit{wr\_way} & 0：写全字；1：带 label；3：不带 label \\
\end{tabular}

\textbf{描述}

地址自增模式适用于连续地址的访问，首次使用需设置 ST-OFFSET 指令来配置起始地址偏移，然后通过 ST-AUTO-DATA 指令将 \textit{Rsrc} 中保留的 16 位数存储至内存地址 \textit{Rdst} + \textit{Offset} 中，存储方式见表 \ref{table:data_format-with-write_way-auto}。其中 ST-AUTO-DATA 的执行次数用 write\_cnt 表示。 \\


\begin{table}[!ht]
    \centering
    \begin{tabular}{ c l l l }
    \textbf{wr\_way}  & \textbf{write\_cnt} & \textbf{存储数据} & \textbf{运算}  \\ \hline
    0 &*&Mem [Rdst + Offset]\{31:0\} =\{PC[10:0], 3'b0, Label[1:0], Rsrc[15:0]\}  & 写全字，包含指针和数据 \\
    1 & 奇数 & Mem [Rdst + Offset]\{15:0\}  = \{Label[1:0],Rscr[13:0]\}& 低半字存储带 label 的数据\\
    1 & 偶数 & Mem [Rdst + Offset]\{31:16\} = \{Label[1:0],Rscr[13:0]\}& 高半字存储带 label 的数据\\
    3 & 奇数 & Mem [Rdst + Offset]\{15:0\}  = Rscr[15:0]& 低半字存储不带 label 的数据  \\
    3 & 偶数 & Mem [Rdst + Offset]\{31:16\} = Rscr[15:0]& 高半字存储不带 label 的数据  \\
    \end{tabular}
    \caption{数据存储格式-地址自增模式}
    \label{table:data_format-with-write_way-auto}
\end{table}

写全字的方式如下：

\regfield{PC 信息}{11}{21}{{}}%
\regfield{}{3}{18}{{0}}%
\regfield{label}{2}{16}{{}}%
\regfield{Rsrc 内容}{16}{0}{{}}%
\begin{figure*}[!ht ]
    \centering
    \caption{MEM[Rdst + Offset] 写全字}
    \label{figure:RTCMEM}
\end{figure*}

\vspace{-2em}
\begin{tabular}{ l l }
\textbf{位} & \textbf{描述} - 见图 \ref{figure:RTCMEM} \\
\textit{[15:0] } & 存储 Rsrc 的内容\\
\textit{[17:16] } & 存储数据标志，用户定义的 2 位无符号数\\
\textit{[20:18]} & 默认 3'b0 \\
\textit{[31:21]} & 存储当前指令的 PC 信息，单位为 32 位字\\
\end{tabular}

\textbf{注意：}
\begin{itemize}
    \item 当存储操作为全字时，每执行完一次 ST-AUTO-DATA，Offset 就会自动加 1。
    \item 当存储操作为半字时，每执行完两次 ST-AUTO-DATA，Offset 就会自动加 1，并且存储顺序为先写低半字，再写高半字。
    \item 该指令仅能以 32 位字为单位进行访问。
    \item Mem 写入的是 RTC\_SLOW\_MEM 慢速内存，ULP 协处理器的地址 0 即相当于主 CPU 的地址 \hexprefix{}50000000。
\end{itemize}

\textcolor{black}{\bfseries 地址指定模式}

\regfield{}{4}{28}{6}%
\regfield{}{3}{25}{{4}}%
\regfield{}{4}{21}{{}}%
\regfield{offset}{11}{10}{{}}%
\regfield{}{1}{9}{{}}%
\regfield{wr\_way}{2}{7}{{}}%
\regfield{upper}{1}{6}{{}}%
\regfield{label}{2}{4}{{}}%
\regfield{Rsrc}{2}{2}{{}}%
\regfield{Rdst}{2}{0}{{}}%
\begin{figure*}[!ht]
    \centering
    \caption{指令类型 - 指定地址模式的数据存储}
    \label{figure:STM}
\end{figure*}


\begin{tabular}{ l l }
\textbf{Operand} & \textbf{描述} - 见图 \ref{figure:STM} \\
\textit{Rdst} & 寄存器 R[0-3]，目标寄存器地址，单位为 32 位字\\
\textit{Rsrc} & 寄存器 R[0-3]，保留需存储的 16 位数 \\
\textit{label} & 数据标志，用户自定义的 2 位无符号数\\
\textit{upper} & 0：写低半字；1：写高半字 \\
\textit{wr\_way} & 0：写全字；1：带 label；3：不带 label \\
\textit{offset} & 11 位有符号数，单位为 32 位字\\
\end{tabular}

\textbf{描述}

指定地址模式主要运用于地址不连续的存储需求，每条指令均需要提供存储地址及偏移量。存储方式见表 \ref{table:data_format-with-write_way}。\\


\begin{table}[!ht]
    \centering
    \begin{tabular}{ c l l l }
    \textbf{wr\_way}  & \textbf{upper} & \textbf{存储数据} & \textbf{运算}  \\ \hline
    0 &*&Mem [ Rdst + Offset ]\{31:0\} =\{PC[10:0], 3'b0, Label[1:0], Rsrc[15:0]\}  & 写全字，包含指针和数据 \\
    1 & 0 & Mem [ Rdst + Offset ]\{15:0\}  = \{Label[1:0],Rscr[13:0]\} & 低半字存储带 label 的数据\\
    1 & 1 & Mem [ Rdst + Offset ]\{31:16\} = \{Label[1:0],Rscr[13:0]\} & 高半字存储带 label 的数据\\
    3 & 0 & Mem [ Rdst + Offset ]\{15:0\}  =  Rsrc[15:0] &低半字存储不带 label 的数据  \\
    3 & 1 & Mem [ Rdst + Offset ]\{31:16\} = Rsrc[15:0]& 高半字存储不带 label 的数据  \\
    \end{tabular}
    \caption{数据存储格式-地址指定模式}
    \label{table:data_format-with-write_way}
\end{table}

\subsubsection {LD - 从内存加载数据}

\regfield{}{4}{28}{{13}}%
\regfield{rd\_upper}{1}{27}{{}}%
\regfield{}{6}{21}{{}}%
\regfield{offset}{11}{10}{{}}%
\regfield{}{6}{4}{{}}%
\regfield{Rsrc}{2}{2}{{}}%
\regfield{Rdst}{2}{0}{{}}%
\begin{figure*}[!ht]
    \centering
    \caption{指令类型 - LD}
    \label{figure:LD}
\end{figure*}

\begin{tabular}{ l l }
\textbf{Operand} & \textbf{描述} - 见图 \ref{figure:LD} \\
\textit{Rdst} & 寄存器 R[0-3]，目标寄存器\\
\textit{Rsrc} & 寄存器 R[0-3]，保留目标寄存器的地址，单位为 32 位字 \\
\textit{Offset} & 11 位有符号数，单位为 32 位字 \\
\textit{rd\_upper} & 读取数据位置选择： \\
 & 1 - 读取高半字\\
 & 0 - 读取低半字
\end{tabular}

\textbf{描述}

该指令可根据 \textit{rd\_upper} 的配置将内存地址 \textit{Rsrc} + \textit{Offset} 中的高或低半字加载至目标寄存器 \textit{Rdst}：

~
\begin{centering}
    Rdst[15:0] = Mem[Rsrc + Offset] \\
\end{centering}
~\\

\textbf{注意：}
\begin{itemize}
    \item 该指令仅能以 32 位字为单位进行访问。
    \item 加载的 Mem 将存储至 RTC\_SLOW\_MEM 慢速内存中，ULP 协处理器的地址 0 即相当于主 CPU 的地址 \hexprefix{}50000000。
\end{itemize}


\subsubsection {JUMP - 跳转至绝对地址} \label{paragraph:JUMP}

\regfield{}{4}{28}{{8}}%
\regfield{}{2}{26}{{1}}%
\regfield{}{1}{25}{{}}%
\regfield{Type}{3}{22}{{}}%
\regfield{Sel}{1}{21}{{}}%
\regfield{}{8}{13}{{}}%
\regfield{ImmAddr}{11}{2}{{}}%
\regfield{Rdst}{2}{0}{{}}%
\begin{figure*}[!ht]
    \centering
    \caption{指令类型 - JUMP}
    \label{figure:JUMP}
\end{figure*}

\begin{tabular}{ l l }
\textbf{Operand} & \textbf{描述} - 见图 \ref{figure:JUMP} \\
\textit{Rdst} & 寄存器 R[0-3]，储存需跳转至的目标地址 \\
\textit{ImmAddr} & 11 位地址，单位为 32 位字\\
\textit{Sel} & 跳转目标地址来源： \\
 & 0 - \textit{ImmAddr} 存储的地址\\
 & 1 - \textit{Rdst} 存储的地址 \\
\textit{Type} & 跳转类型： \\
 & 0 - 无条件跳转 \\
 & 1 - 有条件跳转，仅当最后一次 ALU 运算设置了零标志位时跳转 \\
 & 2 - 有条件跳转，仅当最后一次 ALU 运算设置了溢出标志位时跳转
\end{tabular}

\textbf{注意：} \\
所有跳转地址均以 32 位字为单位。

\textbf{描述} \\
该指令可以让 ULP-FSM 跳转至特定地址，跳转可以为无条件跳转或有条件跳转。


\subsubsection {JUMPR - 跳转至相对地址（基于 R0 寄存器判断）}

\regfield{}{4}{28}{{8}}%
\regfield{}{2}{26}{{0}}%
\regfield{Step}{8}{18}{{}}%
\regfield{Cond}{2}{16}{{}}%
\regfield{Threshold}{16}{0}{{}}%
\begin{figure*}[!ht]
    \centering
    \caption{指令类型 - JUMPR}
    \label{figure:JUMPR}
\end{figure*}

\begin{tabular}{ l l }
\textbf{Operand} & \textbf{描述} - 见图 \ref{figure:JUMPR} \\
\textit{Threshold} & 跳转条件阈值（见下方 \textit{Cond}）\\
\textit{Cond} & 跳转条件：\\
 & 0 - 如果 R0 < \textit{Threshold}，即跳转 \\
 & 1 - 如果 R0 > \textit{Threshold}，即跳转 \\
 & 2 - 如果 R0 = \textit{Threshold}，即跳转 \\
\textit{Step} & 相对位移量，单位为 32 位字：\\
 & 如果 Step[7] = 0，则 PC = PC + Step[6:0] \\
 & 如果 Step[7] = 1，则 PC = PC - Step[6:0] \\
\end{tabular}

\textbf{注意：} \\
所有跳转地址均以 32 位字为单位。

\textbf{描述} \\
如果跳转条件（即比较 R0 寄存器的值与 \textit{Threshold} 阈值）为真，该指令可以让协处理器跳转至 1 个相对地址。

\newpage
\subsubsection {JUMPS - 跳转至相对地址（基于阶段计数器寄存器判断）}\label{jumpstorelativeaddress}

\regfield{}{4}{28}{{8}}%
\regfield{}{2}{26}{{2}}%
\regfield{Step}{8}{18}{{}}%
\regfield{Cond}{2}{16}{{}}%
\regfield{Threshold}{16}{0}{{}}%
\begin{figure*}[!ht]
    \centering
    \caption{指令类型 - JUMPS}
    \label{figure:JUMPS}
\end{figure*}

\begin{tabular}{ l l }
\textbf{Operand} & \textbf{描述} - 见图 \ref{figure:JUMPS} \\
\textit{Threshold} & 跳转条件阈值（见下方 \textit{Cond}）\\
\textit{Cond} & 跳转条件： \\
 & 1X - 如果 \textit{Stage\_cnt} <= \textit{Threshold}，即跳转 \\
 & 00 - 如果 \textit{Stage\_cnt} < \textit{Threshold}，即跳转 \\
 & 01 - 如果 \textit{Stage\_cnt} >= \textit{Threshold}，即跳转\\
 \textit{Step} & 相对位移量，单位为 32 位字： \\
 & 如果 Step[7] = 0，则 PC = PC + Step[6:0] \\
 & 如果 Step[7] = 1，则 PC = PC - Step[6:0] \\
\end{tabular}


\textbf{注意：}
\begin{itemize}
    \item 有关阶段计数器的相关设置，请见章节 \ref{section:alu-cnt} ALU 阶段计数器。
    \item 所有跳转地址均以 32 位字为单位。
\end{itemize}

\textbf{描述} \\
如果跳转条件（即比较 \textit{Stage\_cnt} 阶段计数器寄存器的值与 \textit{Threshold} 阈值）为真，该指令可以让协处理器跳转至 1 个相对地址。


\subsubsection {HALT - 结束程序} \label{subsection:HALT}

\regfield{}{4}{28}{{11}}%
\regfield{}{28}{0}{{}}%
\begin{figure*}[!ht]
    \centering
    \caption{指令类型 - HALT}
    \label{figure:HALT}
\end{figure*}

\textbf{描述} \\
该指令可以让 ULP-FSM 进入断电模式。

\textbf{注意：} \\
执行该指令后，ULP 协处理器的硬件定时器将开始计时。


\subsubsection {WAKE - 唤醒芯片}

\regfield{}{4}{28}{{9}}%
\regfield{}{2}{26}{{0}}%
\regfield{}{25}{1}{{}}%
\regfield{}{1}{0}{{1'b1}}%
\begin{figure*}[!ht]
    \centering
    \caption{指令类型 - WAKE}
    \label{figure:WAKE}
\end{figure*}

\textbf{描述} \\
该指令可以让 ULP-FSM 向 RTC 控制器发送中断。
\begin{itemize}
    \item 当芯片处于 Deep-sleep 模式时，该指令可唤醒芯片。
    \item 当芯片处于 Deep-sleep 之外的模式时，如果 RTC\_CNTL\_INT\_ENA\_REG 寄存器设置了 RTC\_CNTL\_ULP\_CP\_INT\_ENA 中断位，该指令即触发 RTC 中断。
\end{itemize}



\subsubsection {WAIT - 等待若干个周期}

\regfield{}{4}{28}{{4}}%
\regfield{}{12}{16}{{}}%
\regfield{Cycles}{16}{0}{{}}%
\begin{figure*}[!ht]
    \centering
    \caption{指令类型 - WAIT}
    \label{figure:WAIT}
\end{figure*}

\begin{tabular}{ l l }
\textbf{Operand} & \textbf{描述} - 见图 \ref{figure:WAIT} \\
\textit{Cycles} & 等待周期
\end{tabular}

\textbf{描述} \\
该指令可以设定 ULP-FSM 暂停工作的等待周期。


\subsubsection {TSENS - 对温度传感器进行测量}

\regfield{}{4}{28}{{10}}%
\regfield{}{12}{16}{{}}%
\regfield{Wait\_Delay}{14}{2}{{}}%
\regfield{Rdst}{2}{0}{{}}%
\begin{figure*}[!ht]
    \centering
    \caption{指令类型 - TSENS}
    \label{figure:TSENS}
\end{figure*}

\begin{tabular}{ l l }
\textbf{Operand} & \textbf{描述} - 见图 \ref{figure:TSENS} \\
\textit{Rdst} & 目标地址寄存器 R[0-3]，存储测量结果 \\
\textit{Wait\_Delay} & 测量进行的周期数
\end{tabular}

\textbf{描述} \\
增加测量周期数 \textit{Wait\_Delay} 有助于提高测量精确度或优化测量结果。该指令可对片上温度传感器的数据进行测量，并将测量结果存入 1 个通用寄存器中。


\subsubsection {ADC - 对 ADC 进行测量} \label{subsection:ADC}

\regfield{}{4}{28}{{5}}%
\regfield{}{21}{7}{{}}%
\regfield{Sel}{1}{6}{{}}%
\regfield{Sar\_Mux}{4}{2}{{}}%
\regfield{Rdst}{2}{0}{{}}%
\begin{figure*}[!ht]
    \centering
    \caption{指令类型 - ADC}
    \label{figure:ADC}
\end{figure*}

\begin{tabular}{ l l }
\textbf{Operand} & \textbf{描述} - 见图 \ref{figure:ADC} \\
\textit{Rdst} & 目标地址寄存器 R[0-3]，将存储测量结果。\\
\textit{Sar\_Mux} & 使能 SAR ADC 管脚 [Sar\_Mux - 1]，见表 \ref{table:adc inputs1}。 \\
\textit{Sel} & 选择 ADC。0：选择 SAR ADC1；1：选择 SAR ADC2，具体可见表 \ref{table:adc inputs1}。
\end{tabular}

\begin{longtable}{ |p{3.5cm}|c|p{4cm}| }
    \caption{ADC 指令的输入信号}
    \label{table:adc inputs1}
    \\\hline
    \rowcolor{lightgray}
    管脚名 / 信号名 / GPIO \ & \textit{Sar\_Mux} & ADC 选择（\textit{Sel}） \\\hline
    \endfirsthead
    \hline
    \rowcolor{lightgray}
    管脚名 / 信号名 / GPIO \ & \textit{Sar\_Mux} & ADC 选择（\textit{Sel}） \\\hline
    \endhead
    \hline
    \endfoot
    GPIO1    &  1  & \multirow{10}{*}{\textit{Sel} = 0，选择 SAR ADC1} \\ \cline{1-2}
    GPIO2  &  2  & \\ \cline{1-2}
    GPIO3  &  3  & \\ \cline{1-2}
    GPIO4  &  4  & \\ \cline{1-2}
    GPIO5  &  5  & \\ \cline{1-2}
    GPIO6  &  6  & \\ \cline{1-2}
    GPIO7  &  7  & \\ \cline{1-2}
    GPIO8  &  8  & \\ \cline{1-2}
    GPIO9  &  9  & \\ \cline{1-2}
    GPIO10  &  10  & \\ \hline
    GPIO11  &  1  & \multirow{10}{*}{\textit{Sel} = 1，选择 SAR ADC2 } \\ \cline{1-2}
    GPIO12  &  2  & \\ \cline{1-2}
    GPIO13  &  3  & \\ \cline{1-2}
    GPIO14  &  4  & \\ \cline{1-2}
    XTAL\_32k\_P  &  5  & \\ \cline{1-2}
    XTAL\_32k\_N  &  6  & \\ \cline{1-2}
    DAC1  &  7  & \\ \cline{1-2}
    DAC2  &  8  & \\ \cline{1-2}
    GPIO19  &  9  & \\ \cline{1-2}
    GPIO20                 &  10  &
\end{longtable}



\subsubsection {REG\_RD - 从外围寄存器读取} \label{subsection:REG_RD}

\regfield{}{4}{28}{{2}}%
\regfield{High}{5}{23}{{}}%
\regfield{Low}{5}{18}{{}}%
\regfield{}{8}{10}{{}}%
\regfield{Addr}{10}{0}{{}}%
\begin{figure*}[!ht]
    \centering
    \caption{指令类型 - REG\_RD}
    \label{figure:REG_RD}
\end{figure*}

\begin{tabular}{ l l }
\textbf{Operand} & \textbf{描述} - 见图 \ref{figure:REG_RD} \\
\textit{Addr} & 外围设备寄存器地址，单位为 32 位字 \\
\textit{Low} & 寄存器开始位 \\
\textit{High} & 寄存器结束位
\end{tabular}

\textbf{描述}

该指令可以从外设寄存器中读取最高 16 位的内容，并存入通用寄存器。

~
\begin{centering}
    R0 = REG[Addr][High:Low] \\
\end{centering}
~\\
如需读取的内容超过 16 位，即 \textit{High} - \textit{Low} + 1 > 16，则该指令将返回 [Low+15:Low] 的内容。

\textbf{注意：}
\begin{itemize}
    \item 该指令可访问 RTC\_CNTL、RTC\_IO、SENS 及 RTC\_I2C 外围设备中的寄存器。ULP 协处理器可通过相同寄存器在外设总线上的地址 (addr\_peribus1)，计算外围寄存器的地址，具体方式见下：

    ~
    \begin{centering}
        addr\_ulp = (addr\_peribus1 - DR\_REG\_RTCCNTL\_BASE) / 4 \\
    \end{centering}
    ~

    \item \textit{addr\_ulp} 以 32 位字（而非字节）为单位，0 可投射至 DR\_REG\_RTCCNTL\_BASE（从主 CPU 的角度）。因此，10 位 ULP 协处理器的地址可覆盖外围寄存器空间的 4096 字节，包括 DR\_REG\_RTCCNTL\_BASE、DR\_REG\_RTCIO\_BASE、DR\_REG\_SENS\_BASE 及 DR\_REG\_RTC\_I2C\_BASE 区域。
\end{itemize}



\subsubsection {REG\_WR - 写入外围寄存器} \label{subsection:REG_WR}

\regfield{}{4}{28}{{1}}%
\regfield{High}{5}{23}{{}}%
\regfield{Low}{5}{18}{{}}%
\regfield{Data}{8}{10}{{}}%
\regfield{Addr}{10}{0}{{}}%
\begin{figure*}[!ht]
    \centering
    \caption{指令类型 - REG\_WR}
    \label{figure:REG_WR}
\end{figure*}

\begin{tabular}{ l l }
\textbf{Operand} & \textbf{描述} - 见图 \ref{figure:REG_WR} \\
\textit{Addr} & 目标寄存器地址，单位为 32 位字 \\
\textit{Data} & 需写入的值，8 位数 \\
\textit{Low} & 寄存器开始位 \\
\textit{High} & 寄存器结束位
\end{tabular}

\textbf{描述}

该指令最高可以向外设寄存器写入一个 8 位立即值 (Data)。

~
\begin{centering}
    REG[Addr][High:Low] = Data \\
\end{centering}
~\\
如需写入的内容超过 8 位，即 \textit{High} - \textit{Low} + 1 > 8，则该指令会给 8 位以上的内容填充 0。

\textbf{注意：}

有关 \textit{addr\_ulp} 的内容，请见第 \ref{subsection:REG_RD} 节。



\section{ULP-RISC-V} \label{subsection:PicoRV32}
\subsection{ULP-RISC-V 主要特性}
\begin{itemize}
    \item 支持 \href{https://riscv.org/specifications/}{RV32IMC} 指令集
    \item 32 个 32 位通用寄存器
    \item 32 位乘除法器
    \item 支持中断
\end{itemize}

\subsection{乘除法器}
ULP-RISC-V 带有独立的乘除法单元，乘除法相关指令的效率如下表所示：

\begin{longtable}{
  | >{\raggedright\arraybackslash}p{1cm}
  | >{\raggedright\arraybackslash}p{2cm}
  | >{\centering\arraybackslash}p{2.5cm}
  | >{\raggedright\arraybackslash}p{8.5cm}| }
\caption{乘除法指令效率}
\label{table:multi_and_div_instr_zh}
\\\hline
\rowcolor{lightgray}
算法  & 指令 & 执行周期 & 指令描述 \\\hline
\endfirsthead
\hline
\rowcolor{lightgray}
算法  & 指令 & 执行周期 & 指令描述 \\\hline
\endhead
\hline
\endfoot
\hline
\endlastfoot

\multirow{4}{*}{乘}
 & MUL     & 34 & 两个 32 位整数相乘，返回结果低 32 位 \\ \cline{2-4}
 & MULH    & 66 & 两个 32 位的有符号整数相乘，返回结果高 32 位 \\ \cline{2-4}
 & MULHU   & 66 & 两个 32 位的无符号整数相乘，返回结果高 32 位 \\ \cline{2-4}
 & MULHSU  & 66 & 32 位有符号整数与无符号整数相乘，返回结果高 32 位 \\ \hline

\multirow{4}{*}{除}
 & DIV     & 34 & 两个 32 位整数相除，返回商 \\ \cline{2-4}
 & DIVU    & 34 & 两个 32 位的无符号整数相除，返回商 \\ \cline{2-4}
 & REM     & 34 & 两个 32 位的有符号整数相除，返回余数 \\ \cline{2-4}
 & REMU    & 34 & 两个 32 位的无符号整数相除，返回余数 \\
\end{longtable}


\subsection{ULP-RISC-V 中断}
ESP32-S2 的部分传感器、软件以及 RTC I2C 的中断可接入 ULP-RISC-V。用户可配置寄存器 \hyperref[regdesc:SENSSARCOCPUINTENAREG]{SENS\_SAR\_COCPU\_INT\_ENA\_REG} 使能中断，如表 \ref{table:PicoRV32 interrupt list} 所示：

\begin{table}[!ht]
    \centering
    \begin{tabular}{ c l l l }
    \textbf{中断使能位} & \textbf{中断名称} & \textbf{中断描述} \\ \hline
    0 & \label{interrupt:SENSCOCPUTOUCHDONEINT}TOUCH\_DONE\_INT & 触摸传感器完成一个通道扫描则触发此中断 \\
    1 & \label{interrupt:SENSCOCPUTOUCHINACTIVEINT}TOUCH\_INACTIVE\_INT & 触摸传感器的触碰被释放则触发此中断 \\
    2 & \label{interrupt:SENSCOCPUTOUCHACTIVEINT}TOUCH\_ACTIVE\_INT & 触摸传感器感受到触碰则触发此中断 \\
    3 & \label{interrupt:SENSCOCPUSARADC1DONEINT}SARADC1\_DONE\_INT & SAR ADC1 完成一次转换则触发此中断 \\
    4 & \label{interrupt:SENSCOCPUSARADC2DONEINT}SARADC2\_DONE\_INT & SAR ADC2 完成一次转换则触发此中断 \\
    5 & \label{interrupt:SENSCOCPUTSENSINT}TSENS\_DONE\_INT & 温度传感器数据转存完成则触发此中断 \\
    6 & \label{interrupt:SENSCOCPUSTARTINT}RISCV\_START\_INT & ULP-RISC-V 上电开始工作则触发此中断 \\
    7 & \label{interrupt:SENSCOCPUSWINT}SW\_INT & 软件中断 \\
    8 & \label{interrupt:SENSCOCPUSWDINT}SWD\_INT & 超级看门狗超时则触发此中断 \\
    \end{tabular}
    \caption{ULP-RISC-V 的中断列表}
    \label{table:PicoRV32 interrupt list}
\end{table}

\textbf{注意：}
\begin{itemize}
    \item 除了上述中断以外，ULP-RISC-V 还可以响应来自 RTC\_IO 的中断，只需要将 RTC\_IO 配置为输入模式即可。具体触发方式可通过 \hyperref[fielddesc:RTCIOGPIOPINNINTTYPE]{RTCIO\_GPIO\_PIN\regindex{n}\_INT\_TYPE} 进行配置，但只能选择电平类触发。关于 RTC\_IO 的配置，见 IO MUX 和 GPIO 交换矩阵章节。
    \item RTC\_IO 的中断需要通过查询寄存器 RTCIO\_RTC\_GPIO\_STATUS\_INT 识别中断来源，停用 RTC\_IO 可清除此中断。
    \item 软件中断通过配置寄存器 \hyperref[fielddesc:RTCCNTLCOCPUSWINTTRIGGER]{RTC\_CNTL\_COCPU\_SW\_INT\_TRIGGER} 产生。
    \item RTC I2C 的中断描述见章节 \ref{ulp_interrupt}。
    \item 所有中断服务程序入口为 \hexprefix{}10，复位入口为 \hexprefix{}0。
\end{itemize}

当中断发生时，ULP-RISC-V 将执行以下操作：

\begin{enumerate}
    \item 将当前 PC 保存到 Q0
    \item 将被响应的中断保存到 Q1
    \item 跳转到中断入口 (\hexprefix{}10)
\end{enumerate}

进入中断入口后，需要保存上下文，可使用 Q2 和 Q3 作为暂存备份寄存器。当需要退出中断时，恢复上下文，执行 retirq 指令即可从中断返回，跳转到 Q0。ULP 不支持中断嵌套，在执行 retirq 指令之前不会被任何中断打断。

\section{RTC I2C 控制器} \label{subsection:RTC_I2C-Controller}

ULP 协处理器可通过 RTC I2C 控制器与外部 I2C 从机进行基本的读写操作。

\subsection{连接 RTC I2C 信号}

SDA 和 SCL 时钟信号可通过 \hyperref[regdesc:RTCIOSARI2CIOREG]{RTCIO\_SAR\_I2C\_IO\_REG} 寄存器，连接至 2 个 GPIO 管脚（4 个可选），详细定义请见章节 IO MUX 和 GPIO 矩阵中的 RTC\_MUX 管脚清单。

\subsection{配置 RTC I2C 控制器}

ULP 协处理器在正常使用 I2C 指令之前必须配置 RTC I2C 控制器中的特定参数，具体即向 RTC I2C 寄存器写入特定计时参数。这一步可通过主 CPU 或 ULP 自身运行程序完成。

\begin{enumerate}
    \item 通过 \hyperref[fielddesc:RTCI2CSCLLOWPERIODREG]{RTC\_I2C\_SCL\_LOW\_PERIOD\_REG} 和 \hyperref[fielddesc:RTCI2CSCLHIGHPERIODREG]{RTC\_I2C\_SCL\_HIGH\_PERIOD\_REG} 设置 RTC\_FAST\_CLK 周期中 SCL 时钟的高低电平宽度和周期（例，频率为 100 kHz时，设置 \hyperref[fielddesc:RTCI2CSCLLOWPERIODREG]{RTC\_I2C\_SCL\_LOW\_PERIOD\_REG} = 40、 \hyperref[fielddesc:RTCI2CSCLHIGHPERIODREG]{RTC\_I2C\_SCL\_HIGH\_PERIOD\_REG} = 40）。

    \item 通过 RTC\_FAST\_CLK 中的 \hyperref[regdesc:RTCI2CSDADUTYREG]{RTC\_I2C\_SDA\_DUTY\_REG} 设置 SDA 切换前等待的周期数（例，\hyperref[regdesc:RTCI2CSDADUTYREG]{RTC\_I2C\_SDA\_DUTY\_REG} = 16）。

    \item 通过 \hyperref[regdesc:RTCI2CSCLSTARTPERIODREG]{RTC\_I2C\_SCL\_START\_PERIOD\_REG} 设置启动信号后的等待时间（例，\hyperref[regdesc:RTCI2CSCLSTARTPERIODREG]{RTC\_I2C\_SCL\_START\_PERIOD\_REG} = 30）。

    \item 通过 \hyperref[regdesc:RTCI2CSCLSTOPPERIODREG]{RTC\_I2C\_SCL\_STOP\_PERIOD\_REG} 设置停止信号前的等待时间（例，\hyperref[regdesc:RTCI2CSCLSTOPPERIODREG]{RTC\_I2C\_SCL\_STOP\_PERIOD\_REG} = 44）。

    \item 通过 \hyperref[fielddesc:RTCI2CTIMEOUTREG]{RTC\_I2C\_TIME\_OUT\_REG} 设置通信超时参数（例，\hyperref[fielddesc:RTCI2CTIMEOUTREG]{RTC\_I2C\_TIME\_OUT\_REG} = 200）。

    \item 通过 \hyperref[regdesc:RTCI2CCTRLREG]{RTC\_I2C\_CTRL\_REG} 中的 \hyperref[fielddesc:RTCI2CMSMODE]{RTC\_I2C\_MS\_MODE} 位启动主机模式。

    \item 配置外部从机地址：
    \begin{itemize}
        \item 如果选择使用主 CPU 或 ULP-RISC-V，则需要将外部从机的地址写入 \hyperref[fielddesc:SENSSARI2CSLAVEADDR]{SENS\_SAR\_I2C\_SLAVE\_ADDR}。
        \item 如果选择使用 ULP-FSM，则需要将外部从机的地址写入 SENS\_I2C\_SLAVE\_ADDR\regindex{{\it{n}}} (\regindex{{\it{n}}}: 0-7)。
    \end{itemize}
    最多可通过这种方式预编程 8 个从机地址。此后，可为每次通信选择 1 个上述地址，共同组成协处理器 I2C 指令。
\end{enumerate}

完成上述 RTC I2C 初始化配置后，即可由主 CPU 或者协处理器开始与外部 I2C 发起通信。


\subsection{使用 RTC I2C}




\subsubsection{I2C 指令编码格式}

RTC I2C 的指令编码与 I2C0/I2C1 的编码方式保持一致（见 I2C 控制器章节的 CMD\_Controller）。不同的地方在于，RTC I2C 为不同的操作设置了固定的指令段，具体如下：
\begin{itemize}
    \item 命令 0 \verb+~+ 命令 1：I2C 写操作
    \item 命令 2 \verb+~+ 命令 6：I2C 读操作
\end{itemize}

\textbf{注意：}
所有从机地址均为 7 位。


\subsubsection{I2C\_RD - I2C 读取流程}

执行 I2C 读取操作之前，需配置以下信息：

\begin{itemize}
    \item 根据需求配置 I2C 的指令列表，包括指令顺序、指令码和读取数据个数 (byte\_num) 等信息。配置方法见 I2C 控制器章节中 I2C0/I2C1 的配置方法。
    \item 使用寄存器 \hyperref[fielddesc:SENSSARI2CREGADDR]{SENS\_SAR\_I2C\_REG\_ADDR} 配置从机寄存器地址。
    \item 置位 \hyperref[fielddesc:SENSSARI2CSTARTFORCE]{SENS\_SAR\_I2C\_START\_FORCE}，以及 \hyperref[fielddesc:SENSSARI2CSTART]{SENS\_SAR\_I2C\_START} 开始 I2C 的传输。
    \item 每接收到 RTC\_I2C\_RX\_DATA\_INT 中断，将读取的数据 (RTC\_I2C\_RDATA) 转存至 SRAM RTC 慢速内存中或直接使用。
\end{itemize}

I2C\_RD 指令的执行步骤如下，见图 \ref{figure:I2C-read-operation}：

\begin{enumerate}
    \item 主机发送启动 (START) 信号；
    \item 主机发送命令字节，包括从机地址和读/写控制位（此时，读/写控制位置为 0，代表“写”）。从机地址可从 SENS\_I2C\_SLAVE\_ADDR\regindex{{\it{n}}} 中获取；
    \item 从机发送应答 (ACK) 信号；
    \item 主机发送从机寄存器地址；
    \item 从机发送应答信号；
    \item 主机发送重复启动 (RSTART) 信号；
    \item 主机发送从机地址，其中读/写控制位置为 1，代表“读”；
    \item 从机发送 1 个字节的数据；
    \item 主机判断传输字节数是否达到当前指令设定的传输字节数。如果达到规定字数，则从读指令中跳出，主机发送非应答信号。否则重复步骤8，继续等待从机发送下一个字节。
    \item 进入停止指令，主机发送停止 (STOP) 信号，结束读取。
\end{enumerate}

\begin{figure*}[!ht]
    \centering
    \begin{bytefield}[bitwidth=1em]{}
        \bitbox[]{5}{} & \bitbox[]{2}{1} & \bitbox[]{8}{2} & \bitbox[]{1}{3} & \bitbox[]{8}{4} & \bitbox[]{1}{5} & \bitbox[]{2}{6} & \bitbox[]{8}{7} & \bitbox[]{8}{8} & \bitbox[]{1}{9} & \bitbox[]{2}{10} \\
        \bitbox[]{5}{\textbf{Master}} & \bitbox{2}{\rotatebox{90}{\tiny START}} & \bitbox{8}{Slave Address W} & \bitbox{1}{} & \bitbox{8}{Reg Address} & \bitbox{1}{} & & \bitbox{2}{\rotatebox{90}{\tiny RSTRT}} \bitbox{8}{Slave Address R} & \bitbox{8}{} & \bitbox{1}{\rotatebox{90}{\tiny NACK}} & \bitbox{2}{\rotatebox{90}{\tiny STOP}} \\
        \bitbox[]{5}{\textbf{Slave}} & \bitbox{2}{} & \bitbox{8}{} & \bitbox{1}{\rotatebox{90}{\tiny ACK}} & \bitbox{8}{} & \bitbox{1}{\rotatebox{90}{\tiny ACK}} & \bitbox{2}{} & \bitbox{8}{} & \bitbox{8}{Data(n)} & \bitbox{1}{} & \bitbox{2}{}
    \end{bytefield}
    \caption{I2C 读操作}
    \label{figure:I2C-read-operation}
\end{figure*}

\textbf{注意：} \\
RTC I2C 控制器外围设备会对 SCL 时钟下降沿上的 SDA 信号进行采样。如果从机的 SDA 信号在约 0.38 ms 内发生改变，主机则将接收到不正确的数据。


\subsubsection {I2C\_WR - I2C 写流程}

执行 I2C 写操作之前，需配置以下信息：

\begin{itemize}
    \item 根据需求配置 I2C 的指令列表，包括指令顺序、指令码和待写的数据个数 (byte\_num) 等信息。配置方法见 I2C 控制器章节中 I2C0/I2C1 的配置方法；
    \item 使用寄存器 \hyperref[fielddesc:SENSSARI2CREGADDR]{SENS\_SAR\_I2C\_REG\_ADDR} 配置从机寄存器地址；
    \item 使用寄存器 hyperref[fielddesc:SENSSARI2CWDATA]{SENS\_SAR\_I2C\_WDATA} 配置传输数据；
    \item 置位 \hyperref[fielddesc:SENSSARI2CSTARTFORCE]{SENS\_SAR\_I2C\_START\_FORCE}，以及 \hyperref[fielddesc:SENSSARI2CSTART]{SENS\_SAR\_I2C\_START} 开始 I2C 的传输；
    \item 每次当接收到 RTC\_I2C\_TX\_DATA\_INT 中断，更新下一个需要传输的数据 (\hyperref[fielddesc:SENSSARI2CWDATA]{SENS\_SAR\_I2C\_WDATA})。
\end{itemize}


I2C\_WR 指令的执行步骤如下，见图 \ref{figure:I2C-write-operation}：

\begin{enumerate}
    \item 主机发送开始信号；
    \item 主机发送命令字节，包括从机地址和读/写控制位（此时，读/写控制位置为 0，代表“写”）。从机地址可从 SENS\_I2C\_SLAVE\_ADDR\regindex{{\it{n}}} 中获取；
    \item 从机发送应答信号；
    \item 进入下一条指令，主机发送从机寄存器地址；
    \item 从机发送应答信号；
    \item 主机发送重复启动信号；
    \item 主机发送从机地址，其中读/写位置为 0，代表“写”；
    \item 主机发送 1 个字节的数据；
    \item 从机发送应答信号；主机判断传输字节数是否达到当前指令设定的传输字节数，如果达到规定字数，则结束写指令跳转至下一条指令。否则重复步骤 8，继续发送下一个字节；
    \item 进入停止指令，主机发送停止指令，结束本次传输。
\end{enumerate}


\begin{figure*}[!ht]
    \centering
    \begin{bytefield}[bitwidth=1em]{}
        \bitbox[]{5}{} & \bitbox[]{2}{1} & \bitbox[]{8}{2} & \bitbox[]{1}{3} & \bitbox[]{8}{4} & \bitbox[]{1}{5} & \bitbox[]{2}{6} & \bitbox[]{8}{7} & \bitbox[]{8}{8} & \bitbox[]{1}{9} & \bitbox[]{2}{10} \\
        \bitbox[]{5}{\textbf{Master}} & \bitbox{2}{\rotatebox{90}{\tiny START}} & \bitbox{8}{Slave Address W} & \bitbox{1}{} & \bitbox{8}{Reg Address} & \bitbox{1}{} & & \bitbox{2}{\rotatebox{90}{\tiny RSTRT}} \bitbox{8}{Slave Address W} & \bitbox{8}{Data(n)} & \bitbox{1}{} & \bitbox{2}{\rotatebox{90}{\tiny STOP}} \\
        \bitbox[]{5}{\textbf{Slave}} & \bitbox{2}{} & \bitbox{8}{} & \bitbox{1}{\rotatebox{90}{\tiny ACK}} & \bitbox{8}{} & \bitbox{1}{\rotatebox{90}{\tiny ACK}} & \bitbox{2}{} & \bitbox{8}{} & \bitbox{8}{} & \bitbox{1}{\rotatebox{90}{\tiny ACK}} & \bitbox{2}{}
    \end{bytefield}
    \caption{I2C 写操作}
    \label{figure:I2C-write-operation}
\end{figure*}

\subsubsection {检测错误条件} \label{paragraph:ULP-detecting-errors}
应用程序可以通过查询 \hyperref[regdesc:RTCI2CINTSTREG]{RTC\_I2C\_INT\_ST\_REG} 寄存器中的特定位，判断指令是否成功执行。为了检查特定的通信活动，应首先设置 \hyperref[regdesc:RTCI2CINTENAREG]{RTC\_I2C\_INT\_ENA\_REG} 寄存器中的相应位。注意，系统位图将移 1。如果检测到特定通信活动，且设置了 \hyperref[regdesc:RTCI2CINTSTREG]{RTC\_I2C\_INT\_ST\_REG} 寄存器，则可通过 \hyperref[regdesc:RTCI2CINTCLRREG]{RTC\_I2C\_INT\_CLR\_REG} 寄存器清零。

\subsection{RTC I2C 中断}\label{ulp_interrupt}

\begin{itemize}
    \item \label{interrupt:RTCI2CSLAVETRANCOMPINT}RTC\_I2C\_SLAVE\_TRAN\_COMP\_INT：从机完成传输后则触发此中断。
    \item \label{interrupt:RTCI2CARBITRATIONLOSTINT}RTC\_I2C\_ARBITRATION\_LOST\_INT：主机失去总线控制权后则触发此中断。
    \item \label{interrupt:RTCI2CMASTERTRANCOMPINT}RTC\_I2C\_MASTER\_TRAN\_COMP\_INT：主机传输完成后则触发此中断。
    \item \label{interrupt:RTCI2CTRANSCOMPLETEINT}RTC\_I2C\_TRANS\_COMPLETE\_INT：检测到 STOP 位时则触发此中断。
    \item \label{interrupt:RTCI2CTIMEOUTINT}RTC\_I2C\_TIME\_OUT\_INT：出现超时事件则触发此中断。
    \item \label{interrupt:RTCI2CACKERRINT}RTC\_I2C\_ACK\_ERR\_INT：出现 ACK 错误则触发此中断。
    \item \label{interrupt:RTCI2CRXDATAINT}RTC\_I2C\_RX\_DATA\_INT：接收数据则触发此中断。
    \item \label{interrupt:RTCI2CTXDATAINT}RTC\_I2C\_TX\_DATA\_INT：发送数据则触发此中断。
    \item \label{interrupt:RTCI2CDETECTSTARTINT}RTC\_I2C\_DETECT\_START\_INT：检测到开始信号则触发此中断。
\end{itemize}

\section{基地址}\label{base_address}


\subsection{协处理器基地址}

用户可以通过两个不同的寄存器基地址访问协处理器相关的寄存器，如表 \ref{table:base_addr} 所示。更多信息，请访问 \ref{mod:sysmem} \textit{系统与存储器}章节。



\begin{longtable}[c]{|l|l|C{2.5cm}|}
\caption{协处理器基地址}
\label{table:base_addr}
\\\hline
\rowcolor{lightgray}
    模块 &   访问总线      & 基地址\\\hline
    \multirow{2}*{ULP (ALWAYS\_ON)}&  PeriBUS1     & \hexprefix{}3F408000      \\\cline{2-3}
     & PeriBUS2       & \hexprefix{}60008000 \\\hline
    \multirow{2}*{ULP (RTC\_PERI)}&  PeriBUS1     & \hexprefix{}3F408800      \\\cline{2-3}
    & PeriBUS2       & \hexprefix{}60008800 \\\hline
\end{longtable}

其中：
\begin{itemize}
   \item ULP(ALWAYS\_ON) 为非掉电寄存器，寄存器不会随着 RTC\_PERI 的电源域（见章节 \ref{mod:lowpowm} \textit{\nameref{mod:lowpowm}}）掉电而复位。
    \item ULP(RTC\_PERI) 处于 RTC\_PERI 电源域下，如果 RTC\_PERI 的电源域（见章节 \ref{mod:lowpowm} \textit{\nameref{mod:lowpowm}}）掉电，寄存器值将会被复位。
\end{itemize}

\subsection{RTC I2C 基地址}
与 RTC I2C 相关的寄存器有两组：RTC\_PERI 与 I2C，并且都处于 RTC\_PERI 的电源域下。用户都可以通过两个不同的寄存器基地址进行访问，如表 \ref{table:i2c_base_addr} 所示。更多信息，请访问 \ref{mod:sysmem} \textit{系统与存储器}章节。

\begin{longtable}[c]{|l|l|C{2.5cm}|}
\caption{RTC I2C 基地址}
\label{table:i2c_base_addr}
\\\hline
\rowcolor{lightgray}
    模块 &   访问总线      & 基地址\\\hline
    \multirow{2}*{RTC I2C (RTC\_PERI)}&  PeriBUS1     & \hexprefix{}3F408800      \\\cline{2-3}
     & PeriBUS2       & \hexprefix{}60008800 \\\hline
    \multirow{2}*{RTC I2C (I2C)}&  PeriBUS1     & \hexprefix{}3F408C00      \\\cline{2-3}
    & PeriBUS2       & \hexprefix{}60008C00 \\\hline
\end{longtable}

\section{地址映射}
\label{sec:address-mapping}

表 \ref{tab: address-mapping} 列出了 ULP 协处理器访问外设寄存器所用到的基地址寄存器以及相关的地址映射。

\begin{longtable}[c]{ |m{3.0cm}|m{4cm}|c|c|c| }
    \caption{地址映射}
    \label{tab: address-mapping}
    \\\hline
    \rowcolor{lightgray}
    \textbf{外设} & \textbf{基地址寄存器} & \textbf{总线地址} & \textbf{ULP-FSM 基地址} & \textbf{ULP-RISC-V 基地址} \\\hline
    \endfirsthead
    \hline
    \rowcolor{lightgray}
    外设 & 基地址寄存器 & 总线基地址 & ULP-FSM 基地址 & ULP-RISC-V 基地址\\\hline
    \endhead
    \hline
    \endfoot
   RTC Control & DR\_REG\_RTCCNTL\_BASE & \hexprefix{}3F408000 & \hexprefix{}8000 & \hexprefix{}8000 \\\hline
   RTC GPIO & DR\_REG\_RTC\_IO\_BASE & \hexprefix{}3F408400 &\hexprefix{}8400 & \hexprefix{}A400 \\\hline
   ADC, Touch, TSENS & DR\_REG\_SENS\_BASE & \hexprefix{}3F408800 & \hexprefix{}8800& \hexprefix{}C800\\\hline
   RTC I2C & DR\_REG\_RTC\_I2C\_BASE & \hexprefix{}3F408C00 & \hexprefix{}8C00 & \hexprefix{}EC00 \\\hline
\end{longtable}

表 \ref{tab:assiable-registers} 列出了 ULP 协处理器可访问的外设寄存器。

\begin{table}[H]
    \centering
    \caption{ULP 协处理器可访问的外设寄存器}\label{tab:assiable-registers}

    \begin{threeparttable}
    \begin{tabular}[c]{|p{6cm}|p{10cm}|}
    \hline
        \rowcolor{lightgray}
\textbf{外设寄存器} & \textbf{寄存器描述} \\\hline
RTC CNTL 寄存器 & 描述见章节 \ref{mod:lowpowm} \textit{\nameref{mod:lowpowm}} \\\hline
RTC GPIO 寄存器& 描述见章节 \ref{mod:iomuxgpio} \textit{\nameref{mod:iomuxgpio}} \\\hline
ARC、Touch、TSENS 寄存器 & 描述见章节 \ref{mod:sensor} \textit{\nameref{mod:sensor}} \\\hline
RTC I2C 寄存器 & 见本章节第 \ref{ulp-reg-summ} 小节 \textit{\nameref{ulp-reg-summ}} \\\hline
\end{tabular}
     \end{threeparttable}
    \end{table}

\hypertarget{ulp-reg-summ}{}
\section{寄存器列表}
\label{ulp-reg-summ}


请注意，表中的地址为相对于基地址的地址偏移量（相对地址）。请参考章节 \ref{base_address} 获取基地址相关信息。

\subsection{ULP (ALWAYS\_ON) 寄存器列表}
\subfile{\modulefiles/cocpu_always_on_summary_cn.tex}

\subsection{ULP (RTC\_PERI) 寄存器列表}
\subfile{\modulefiles/cocpu_rtc_peri_summary_cn.tex}

\subsection{RTC I2C (RTC\_PERI) 寄存器列表}
\subfile{\modulefiles/i2c_controll_rtc_peri_summary_cn.tex}

\subsection{RTC I2C (I2C) 寄存器列表}
\subfile{\modulefiles/rtc_i2c_reg_summary_cn.tex}

\section{寄存器}

请注意，表中的地址为相对于基地址的地址偏移量（相对地址）。请参考章节 \ref{base_address} 获取基地址相关信息。

\subsection{ULP (ALWAYS\_ON) 寄存器}
\subfile{\modulefiles/cocpu_always_on_cn.tex}

\subsection{ULP (RTC\_PERI) 寄存器}
\subfile{\modulefiles/cocpu_rtc_peri_cn.tex}

\subsection{RTC I2C (RTC\_PERI) 寄存器}
\subfile{\modulefiles/i2c_controll_rtc_peri_cn.tex}

\subsection{RTC I2C (I2C) 寄存器}
\subfile{\modulefiles/rtc_i2c_reg_cn.tex}

\end{document}
